\documentclass{article}

\usepackage{amsmath,amssymb}
\usepackage{xurl}
\usepackage[hidelinks]{hyperref}
\setlength{\emergencystretch}{2em}

% Shared commands for notes in docs/paper-gaps/.

\DeclareUnicodeCharacter{2080}{\ifmmode{}_{0}\else\textsubscript{0}\fi}
\DeclareUnicodeCharacter{2081}{\ifmmode{}_{1}\else\textsubscript{1}\fi}
\DeclareUnicodeCharacter{2082}{\ifmmode{}_{2}\else\textsubscript{2}\fi}
\DeclareUnicodeCharacter{2099}{\ifmmode{}_{n}\else\textsubscript{n}\fi}

\newcommand{\Lean}{\textsc{Lean}}
\ExplSyntaxOn
\NewDocumentCommand{\leanid}{m}{
  \ifmmode
    \textsf{\detokenize{#1}}
  \else
    \group_begin:
    \urlstyle{sf}
    \tl_set:Nn \l_tmpa_tl {#1}
    \tl_replace_all:Nnn \l_tmpa_tl {\_} {_}
    \exp_args:NV \path \l_tmpa_tl
    \group_end:
  \fi
}
\ExplSyntaxOff
\providecommand{\tr}{\operatorname{tr}}
\newcommand{\tnleanIssueBase}{https://github.com/LionSR/TNLean/issues/}
\newcommand{\tnleanPrBase}{https://github.com/LionSR/TNLean/pull/}
\newcommand{\tnleanDocBase}{https://sirui-lu.com/TNLean/docs/}
\newcommand{\ghissue}[1]{\href{\tnleanIssueBase#1}{GitHub issue \##1}}
\newcommand{\ghpr}[1]{\href{\tnleanPrBase#1}{PR \##1}}
\newcommand{\doclink}[1]{\href{\tnleanDocBase#1.html}{\path{#1}}}

% Dirac notation and the identity operator, as in blueprint/src/macros/common.tex.
% \providecommand keeps note-local definitions authoritative.
\usepackage{dsfont}
\providecommand{\ket}[1]{|#1\rangle}
\providecommand{\bra}[1]{\langle #1|}
\providecommand{\braket}[2]{\langle #1 | #2 \rangle}
\providecommand{\ketbra}[2]{|#1\rangle\!\langle #2|}
\providecommand{\Id}{\mathds{1}}
\providecommand{\one}{\mathds{1}}
\providecommand{\C}{\mathbb{C}}
\providecommand{\MN}[1]{M_{#1}(\C)}
\providecommand{\MD}{M_D(\C)}

% External referencing for paper sources.  The build helper writes sanitized
% auxiliary files under build/paper-aux/ and excludes duplicated source labels.
\usepackage{xr}

% Import a generated paper auxiliary file with a caller-supplied label prefix.
% The candidate paths support builds from docs/paper-gaps/, the repository root,
% and build/paper-aux/ itself.
\newcommand{\importpaperaux}[2]{%
  \IfFileExists{../../build/paper-aux/#2.aux}{%
    \externaldocument[#1]{../../build/paper-aux/#2}%
  }{%
    \IfFileExists{build/paper-aux/#2.aux}{%
      \externaldocument[#1]{build/paper-aux/#2}%
    }{%
      \IfFileExists{#2.aux}{%
        \externaldocument[#1]{#2}%
      }{}%
    }%
  }%
}

% General external reference macros
\newcommand{\paperref}[2]{\ref{#1-#2}}
\newcommand{\papereqref}[2]{(\ref{#1-#2})}

% CPSV16 source (1606.00608 / MPDO-22-12-17-2.tex) external document and shortcuts
\importpaperaux{cpsv16-}{1606.00608/MPDO-22-12-17-2-tnlean}
\newcommand{\cpsvref}[1]{\paperref{cpsv16}{#1}}
\newcommand{\cpsveqref}[1]{\papereqref{cpsv16}{#1}}

% Verdict marker: \gapnote{<kind>}{<status>}. Kind is the policy
% classification (clarification, local-correction, scope-restriction,
% unfaithful, false-source, open-gap); status is open, wip (elimination
% underway), resolved, or historical. Machine-read by the site index;
% prints a verdict line.
\newcommand{\gapnote}[2]{\par\noindent\textbf{Verdict:} #1 (#2).\par}


\title{Formalization-Derived Corrections to Wolf's \emph{Quantum Channels \& Operations}}
\author{}
\date{2026-07-13}

\begin{document}
\sloppy
\maketitle

\begin{abstract}
This note records only discrepancies in Michael M.~Wolf's \emph{Quantum
Channels \& Operations: Guided Tour} that were exposed by the present Lean
formalization or by a paper-gap analysis required to state a formal theorem
faithfully. It is not a general errata list for the lecture notes. Its main
body records only formalization-derived corrections; the final appendix records
three separately audited local PDF--transcription errors. Each item quotes the
relevant printed statement, gives the locally correct replacement, and explains
why the correction is necessary.
\end{abstract}

\section{Scope}

The source PDFs are stored in \path{Notes/WolfNotePDF/}; the local
transcription is in \path{Notes/WolfNoteTexSource/}. The corresponding
formalization and paper-gap records are named explicitly below. A distinction
is essential:
\begin{itemize}
  \item a \emph{source error} is a printed mathematical assertion that is false
    as stated;
  \item a \emph{well-posedness correction} makes an implicit boundary convention
    explicit, without changing the substantive theorem;
  \item a \emph{formalization scope gap} is not an erratum: it records that the
    current Lean theorem is narrower than the source or that an auxiliary proof
    remains incomplete.
\end{itemize}

Only the first two kinds belong in the formalization-derived main body. The
appendix keeps transcription errors visibly separate from that record.

\section{The trace sign in the Lorentz-cone converse}
\label{sec:lorentz}

\paragraph{Formalization trigger.}
The formal proof of the converse trace inequality requires nonnegativity of the
trace. Its declaration is
\path{Matrix.IsHermitian.posSemidef_of_trace_re_nonneg_of_card_sub_one_mul_trace_sq_re_le}
in \doclink{TNLean/Analysis/MatrixTraceInequalities}. The discrepancy is
tracked in \path{docs/paper-gaps/wolf_ch3_lorentz_cone_trace_sign.tex}.

\paragraph{Printed source and its role.}
The proposition occurs at the end of the discussion of automorphisms of the
positive cone in Chapter~3.  Its first part gives the elementary necessary
inequality for a positive semidefinite Hermitian matrix; the second part is
intended to characterize the corresponding Lorentz cone.  Chapter~3,
Proposition~3.7 in the local PDF (book p.~66; cited as Proposition~3.9 in a
different numbering used by the formalization) states, for Hermitian $A$,
\begin{equation}\tag{printed converse}
  \operatorname{tr}(A)^2\geq(d-1)\operatorname{tr}(A^2)
  \quad\Longrightarrow\quad A\geq0.
\end{equation}

\paragraph{Correction.}
The valid converse is
\begin{equation}\tag{formalized converse}
  A=A^\dagger,\qquad
  \operatorname{tr}(A)\geq0,\qquad
  \operatorname{tr}(A)^2\geq(d-1)\operatorname{tr}(A^2)
  \quad\Longrightarrow\quad A\geq0.
\end{equation}
In the complex-matrix Lean statement the traces are expressed through real
parts; for Hermitian matrices these are the same real quantities.

\paragraph{Why the correction is necessary.}
Take $A=-I_d$. It is Hermitian and not positive semidefinite, but
\begin{equation}
  \operatorname{tr}(A)^2=d^2
  \geq d(d-1)=(d-1)\operatorname{tr}(A^2).
\end{equation}
Thus the squared inequality alone describes both orientations of the Lorentz
cone. The condition $\operatorname{tr}(A)\geq0$ selects the positive sheet.
It is also used by the proof: when a negative eigenvalue is present, the source
bounds $\operatorname{tr}(A)$ above by a nonnegative quantity and then squares
that bound. Squaring is justified only after fixing the sign of the trace.

\paragraph{Classification.}
This is a genuine source error. The unrestricted printed converse should not
be claimed as formalized; the formal theorem proves the corrected,
trace-nonnegative statement.

\section{The nonempty-family boundary in the Radon--Nikodym theorem}
\label{sec:radon-nikodym}

\paragraph{Formalization trigger.}
The Lean theorem
\leanid{IsCPMap.radon_nikodym_of_stinespring} makes the finite index type
nonempty. The boundary is documented in
\path{docs/paper-gaps/wolf_radon_nikodym_nonempty_family.tex}.

\paragraph{Printed source and its role.}
This theorem is the immediate instrument-valued corollary of the preceding
comparison theorem for ordered completely positive maps.  A supplied
Stinespring representation of the total map is held fixed, and the theorem
resolves each CP component by a positive effect on the same dilation space.
Chapter~2, Theorem~2.4 (book p.~39) says:
\begin{quote}
``Let $\{T_i\}$ be a set of completely positive linear maps such that
$\sum_iT_i=T$ \ldots\ Then there is a set of positive operators $P_i\in M_r$
which sum up to $I_r$ such that
$T_i(A)=V^\dagger(A\otimes P_i)V$.''
\end{quote}

\paragraph{Formal statement.}
The formal theorem uses a \emph{nonempty finite family}
$\{T_i\}_{i\in\iota}$ and concludes
\begin{equation}
  \sum_{i\in\iota}P_i=I_r,
  \qquad
  T_i(A)=V^\dagger(A\otimes P_i)V.
\end{equation}

\paragraph{Why the boundary is needed.}
If the index type is empty, then $\sum_iT_i=T$ forces $T=0$. Taking the zero
Stinespring matrix $V=0$ still satisfies the representation hypothesis, but the
printed conclusion would require
\begin{equation}
  0=\sum_iP_i=I_r,
\end{equation}
which is false when the dilation space has nonzero dimension. For the intended
notion of a quantum instrument, a nonempty finite outcome family is the
substantive case.

\paragraph{Classification.}
This is a local well-posedness correction, not an assertion that the
nonempty-case Radon--Nikodym theorem is false. The source's word ``set'' and
summation notation leave the boundary implicit; Lean must state it.

\section{Suppressed zero multiplicities in the two-pure-state spectrum}
\label{sec:two-pure-state-spectrum}

\paragraph{Formalization trigger.}
The Lean theorem
\leanid{Projectivization.charpoly_add_pureStateMatrix_of_two_le} keeps the full
characteristic polynomial of a sum of two normalized pure-state matrices. The
correction is documented in
\path{docs/paper-gaps/wolf_ch01_two_pure_state_zero_multiplicities.tex}.

\paragraph{Printed source and its role.}
In the spectrum-preserver argument following Wigner's theorem, Chapter~1 says
that the spectrum of two Hermitian rank-one projections $P,Q$ is
\begin{equation}
  1\pm\sqrt{\operatorname{tr}(PQ)}.
\end{equation}
This pair describes the restriction of $P+Q$ to the span of the two pure-state
ranges.

\paragraph{Correction.}
For $d\geq2$, the multiplicity-aware statement is
\begin{equation}
  \chi_{P+Q}(X)
  =X^{d-2}\bigl((X-1)^2-\operatorname{tr}(PQ)\bigr).
\end{equation}
The source shorthand suppresses the $d-2$ zero eigenvalues and their
multiplicities. Dimension one has the separate formula $\chi_{P+Q}(X)=X-2$,
and dimension zero has no projective pure states.

\paragraph{Classification.}
This is a local correction to the displayed spectrum shorthand, not a claim
that Wolf's intended spectrum-preserver argument is false. The two displayed
values contain the transition-probability information used in that argument;
the characteristic polynomial records the complete multiplicity data.

\section{Formalization gaps that are not source corrections}

The following Wolf-specific paper-gap notes are intentionally \emph{not}
errata entries:
\begin{itemize}
  \item \path{wolf_prop28_square_case.tex}: the current normal-form theorem is
    the square-dimensional special case of Wolf's rectangular statement;
  \item \path{wolf_lemma_3_1_top_index_scope.tex},
    \path{wolf_prop_3_2_top_index_scope.tex}, and
    \path{wolf_t_eta_top_index_scope.tex}: endpoint Schmidt-rank cases that
    were formalization scope gaps and are now resolved;
  \item \path{wolf_reduction_criterion_schmidt_premise.tex}: the
    Schmidt-number premise was a formalization dependency and is now resolved;
  \item \path{wolf_ex3_1_choi_positivity_subcase_scope.tex}: the development
    proves selected parameter families, not the full source range;
  \item \path{wolf_ch5_operator_jensen_lieb.tex} and
    \path{wolf_cor67_maximal_support_restriction.tex}: these describe
    formalization boundaries and their resolution, not a false printed theorem.
\end{itemize}

\section{Conclusion}

At present the formalization-derived record contains one genuine mathematical
source error (Section~\ref{sec:lorentz}), one explicit boundary convention
(Section~\ref{sec:radon-nikodym}), and one local correction to a spectrum
shorthand (Section~\ref{sec:two-pure-state-spectrum}). These are the corrections
asserted by the Lean documentation itself.

\appendix
\section{Corrected errors in the local LaTeX transcription}
\label{app:transcription}

This appendix is deliberately placed last.  Its entries were found by comparing
the source PDFs with the local transcription; they are \emph{not} claims that a
Lean proof forced a correction to Wolf's statement.  Each was a mathematically
material transcription error and has now been corrected in the local TeX.

\subsection{The Choi marginal: $\tau_B$ versus $\operatorname{tr}_B\tau$}

\textbf{Source context.}  In Chapter~2, immediately after deriving the
coordinate expression for the Choi--Jamiolkowski correspondence in
Eq.~(2.4), Wolf explains the channel--state dictionary.  For a channel
$T:M_d\to M_{d'}$, the first factor is the output factor $A$ and the second is
the untouched input/reference factor $B$.  The PDF (book p.~34) states:
\begin{quote}
``If $T$ is a quantum channel in the Schr\"odinger picture, then $\tau$ is a
density operator with reduced density matrix $\tau_B=\one/d$.''
\end{quote}
Thus $\tau_B$ denotes $\operatorname{tr}_A\tau$.

\textbf{Transcription and correction.}  A previous local transcription at
\path{ch02_representations.tex:113--114} changed the displayed marginal to
$\operatorname{tr}_B\tau=\one/d$.  It has been corrected to retain Wolf's
notation; equivalently, one may spell out
\begin{equation}
  \tau_B=\operatorname{tr}_A\tau=\frac{\one_d}{d}.
\end{equation}

\textbf{Reason.}  Trace preservation is equivalent to the last identity.  The
other marginal satisfies $\operatorname{tr}_B\tau=T(\one)/d$ and is maximally
mixed precisely when $T$ is unital.  Hence the local transcription accidentally
turns a trace-preserving assertion into a stronger and generally false unitality
assertion.

\subsection{The conjugation in the proof of Kramer's theorem}

\textbf{Source context.}  Chapter~3 proves Kramer's theorem for a Hermitian
$H$ commuting with an anti-unitary symmetry $T=\Gamma V$, where $V^T=-V$.
After showing that $V^\dagger\lvert\overline\psi\rangle$ is another eigenvector,
the PDF (book p.~62, Eq.~(3.29)) establishes orthogonality by the chain
\begin{equation}\tag{PDF (3.29)}
 \langle\psi\mid V^\dagger\mid\overline\psi\rangle
 =-\langle\psi\mid\overline V\mid\overline\psi\rangle
 =-\overline{\langle\overline\psi\mid V\mid\psi\rangle}.
\end{equation}
The final conjugation is then identified with
$\langle\psi\mid V^\dagger\mid\overline\psi\rangle$, proving that this scalar
is its own negative.

\textbf{Transcription and correction.}  A previous local TeX version at
\path{ch03_positive_not_completely.tex:515} omitted the outer bar in the final
term.  It has been corrected to the source-faithful and mathematically
necessary term $-\overline{\langle\overline\psi\mid V\mid\psi\rangle}$.

\textbf{Reason.}  Componentwise,
\begin{equation}
 \langle\psi\mid\overline V\mid\overline\psi\rangle
 =\overline{\langle\overline\psi\mid V\mid\psi\rangle}.
\end{equation}
Without the conjugation the equality is false in general, and the final
self-cancellation in the orthogonality argument does not follow.

\subsection{The order of the conjugated factors in Example~5.3}

\textbf{Source context.}  Chapter~5 studies the adjoint of the map
$T(X)=\tfrac12(X^T+\tfrac12\operatorname{tr}(X)\one)$ as an example of a
Schwarz map that need not be $2$-positive.  In the traceless case, the PDF
(book p.~77, Eq.~(5.10)) subtracts the product
\begin{equation}\tag{PDF (5.10), relevant term}
  \frac14\,\overline A A^T
\end{equation}
from $T^*(A^\dagger A)$.

\textbf{Transcription and correction.}  A previous local TeX version at
\path{ch05_schwarz_inequalities.tex:373,377} instead wrote
$A\overline A^{\,T}$.  It has been corrected to $\overline A A^T$, as in the
source.

\textbf{Reason.}  Since $T^*(A)=A^T/2$ for traceless $A$,
\begin{equation}
 T^*(A^\dagger)T^*(A)
 =\frac14(A^\dagger)^T A^T
 =\frac14\,\overline A A^T.
\end{equation}
The transcribed matrix $A\overline A^{\,T}=AA^\dagger$ is generally a
different operator, even though both matrices are positive semidefinite and
the final positivity conclusion in the example survives.

\end{document}